% Generated. Do not edit.
% Projection scriptural-date-location/evidence, version 1: the Scriptural Date and
% Location page, generated whole from research/chronology.toml, the Scripture
% chronology corpus and the source library by tools/tpt document-types refresh.
% document: liturgy/roman-rite/1962/propers/temporal/46-sixth-after-pentecost
\section*{Appendix: Evidence for the Scriptural Dates}
\label{sec:date-location-evidence}
{\small For each passage, as the Scriptural Date and Location table names it: what is dated, how the corpus weighs each answer, the source's own words, and the sources cited. A subject the corpus does not date is reported as such and given no date.\par}
{\small
\par\smallskip\noindent\textit{Offertory: Ps. 16:5 (Vg.; Heb. 17), Ps. 16:6--7 (Vg.; Heb. 17)}\par\nobreak
\begin{itemize}[leftmargin=1.2em,itemsep=1pt,topsep=1pt,parsep=0pt]
\item \textbf{Historical setting}. No date: the chronology corpus has not yet dated this (research pending).
\item \textbf{Composition} (preferred). Psalms: Before c.~165 B.C. (``before the Maccabean period, around 165 B.C.; no individual psalm can be dated securely''; Catholic critical chronology; New American Bible, Revised Edition --- USCCB web state (2026), Psalms introduction).
\end{itemize}
\par\smallskip\noindent\textit{Communion: Ps. 26:6 (Vg.; Heb. 27)}\par\nobreak
\begin{itemize}[leftmargin=1.2em,itemsep=1pt,topsep=1pt,parsep=0pt]
\item \textbf{Historical setting}. As under Offertory: Ps. 16:5 (Vg.; Heb. 17), Ps. 16:6--7 (Vg.; Heb. 17).
\item \textbf{Composition}. As under Offertory: Ps. 16:5 (Vg.; Heb. 17), Ps. 16:6--7 (Vg.; Heb. 17).
\end{itemize}
\par\smallskip\noindent\textit{Introit: Ps. 27:8--9 (Vg.; Heb. 28), Ps. 27:1 (Vg.; Heb. 28)}\par\nobreak
\begin{itemize}[leftmargin=1.2em,itemsep=1pt,topsep=1pt,parsep=0pt]
\item \textbf{Historical setting}. As under Offertory: Ps. 16:5 (Vg.; Heb. 17), Ps. 16:6--7 (Vg.; Heb. 17).
\item \textbf{Composition}. As under Offertory: Ps. 16:5 (Vg.; Heb. 17), Ps. 16:6--7 (Vg.; Heb. 17).
\end{itemize}
\par\smallskip\noindent\textit{Alleluia: Ps. 30:2--3 (Vg.; Heb. 31)}\par\nobreak
\begin{itemize}[leftmargin=1.2em,itemsep=1pt,topsep=1pt,parsep=0pt]
\item \textbf{Historical setting}. No date: the corpus holds a date for this that it reports from another edition by exception, which is not printed until the corpus records the source's own qualification of it (research pending).
\item \textbf{Historical setting} (preferred). David's flight from the face of Absalom: When he fled from the face of his son Absalom (The Holy Bible, Douay-Rheims (Challoner revision), Project Gutenberg eBook 1581: Ps. 3:1).
\item \textbf{Composition}. As under Offertory: Ps. 16:5 (Vg.; Heb. 17), Ps. 16:6--7 (Vg.; Heb. 17).
\end{itemize}
\par\smallskip\noindent\textit{Gradual: Ps. 89:13 (Vg.; Heb. 90), Ps. 89:1 (Vg.; Heb. 90)}\par\nobreak
\begin{itemize}[leftmargin=1.2em,itemsep=1pt,topsep=1pt,parsep=0pt]
\item \textbf{Historical setting}. As under Offertory: Ps. 16:5 (Vg.; Heb. 17), Ps. 16:6--7 (Vg.; Heb. 17).
\item \textbf{Composition}. As under Offertory: Ps. 16:5 (Vg.; Heb. 17), Ps. 16:6--7 (Vg.; Heb. 17).
\end{itemize}
\par\smallskip\noindent\textit{Gospel: Mark 8:1--9}\par\nobreak
\begin{itemize}[leftmargin=1.2em,itemsep=1pt,topsep=1pt,parsep=0pt]
\item \textbf{Narrated event}. No date: the chronology corpus has not yet dated this (research pending).
\item \textbf{Composition} (disputed). The Gospel of St Mark: A.D.~38--40 (``the tenth or twelfth year after the Ascension (A.D. 38-40)''; The Catholic Encyclopedia, Volume 9 (1910)); A.D.~40 (The Catholic Encyclopedia, Volume 9 (1910)); A.D.~43 (``the third year of Claudius (A.D. 43)''; The Catholic Encyclopedia, Volume 9 (1910)); A.D.~50--67 (``somewhere between A.D. 50 and 67''; The Catholic Encyclopedia, Volume 9 (1910)); A.D.~60--70 (``between 60 and 70''; The Catholic Encyclopedia, Volume 14 (New York, 1912)); After Matthew, and not as late as the destruction of Jerusalem (De auctore, de tempore compositionis et de historica veritate Evangeliorum secundum Marcum et secundum Lucam (1912), responsa-i-ix). Place: Rome traditionally; Egypt reported by Chrysostom (The Catholic Encyclopedia, Volume 9 (1910)).
\end{itemize}
\par\smallskip\noindent\textit{Epistle: Rom. 6:3--11}\par\nobreak
\begin{itemize}[leftmargin=1.2em,itemsep=1pt,topsep=1pt,parsep=0pt]
\item \textbf{Composition} (disputed). The Epistle to the Romans: A.D.~57 (``(Romans), 57''; The Catholic Encyclopedia, Volume 11 (New York, 1911)); A.D.~58 (``the Epistle to the Romans (58)''; The Catholic Encyclopedia, Volume 5 (New York, 1909)).
\end{itemize}
\par}
